\documentclass{article}

\usepackage{booktabs}
\usepackage{tabularx}
\usepackage{hyperref}
\usepackage{float}


\hypersetup{
    colorlinks=true,       % false: boxed links; true: colored links
    linkcolor=red,          % color of internal links (change box color with linkbordercolor)
    citecolor=green,        % color of links to bibliography
    filecolor=magenta,      % color of file links
    urlcolor=cyan           % color of external links
}

\title{Hazard Analysis\\\progname}

\author{\authname}

\date{}

\input{../Comments.text}
\input{../Common.text}

\newcounter{ssrnum} %Safety and Security Requirement (SSR) Number
\newcommand{\ssrthessrnum}{SSR\thessrnum}
\newcommand{\ssrref}[1]{SSR\ref{#1}}

\begin{document}

\maketitle
\thispagestyle{empty}

~\newpage

\pagenumbering{roman}

\begin{table}[hp]
\caption{Revision History} \label{TblRevisionHistory}
\begin{tabularx}{\textwidth}{llX}
\toprule
\textbf{Date} & \textbf{Developer(s)} & \textbf{Change}\\
\midrule
Date1 & Name(s) & Description of changes\\
Date2 & Name(s) & Description of changes\\
... & ... & ...\\
\bottomrule
\end{tabularx}
\end{table}

~\newpage

\tableofcontents

~\newpage

\pagenumbering{arabic}

\wss{You are free to modify this template.}

\section{Introduction}

%\wss{You can include your definition of what a hazard is here.}

A hazard is a system condition or failure that could contribute to an undesirable outcome, resulting in losses 
like inaccurate accreditation assessments, poor deadline recommendations, a loss of course information, or sensitive data breaches.\\

The Course Deadline Management system is intended to support McMaster University's Computing and Software Department 
in reviewing Software Engineering course outlines for compliance with applicable university policies and Canadian 
Engineering Accreditation Board (CEAB) requirements. The system also analyzes course schedules to recommend assignment 
and midterm deadlines that reduce overlapping student workloads, while leaving final scheduling decisions to instructors.\\

This document identifies potential hazards associated with the Course Deadline Management system, examines their consequences, 
and establishes a basis for developing safety and security requirements to mitigate these risks.\\


\section{Scope and Purpose of Hazard Analysis}

%\wss{You should say what \textbf{loss} could be incurred because of the hazards.}

The purpose of this hazard analysis is to identify potential failures in the Course Deadline 
Management system that could negatively affect accreditation compliance reviews, course scheduling, 
or the handling of course information. The analysis covers the processing of course outlines, 
evaluation against McMaster policies and CEAB requirements, deadline conflict detection, and the 
generation of scheduling recommendations. The scope is limited to Software Engineering courses 
within the CAS department and excludes first-year courses, courses outside the department, 
and the security of the hosting infrastructure itself.\\

\noindent Potential losses resulting from system hazards include:

\begin{itemize}
    \item Accreditation non-compliance, resulting in unaddressed violations of McMaster policies or CEAB requirements.
    \item Increased student workload conflicts due to overlapping assignment and midterm deadlines, potentially contributing to academic stress and more MSAF requests.
    \item Incorrect scheduling decisions resulting from inaccurate or misleading recommendations.
    \item Loss or corruption of course information, resulting in inaccurate course outline data and scheduling issues.
    \item Unauthorized disclosure of confidential information, including the master course sheet, compliance findings, and tentative assessment dates.
    \item Wasted administrator and instructor time due to false compliance findings, incorrectly identified deadline conflicts, or inaccurate recommendations.
    \item Loss of trust in the system due to repeated inaccurate or unreliable results, potentially leading users to stop using the system and reducing its intended benefits.
\end{itemize}


\section{System Boundaries and Components}

\wss{Dividing the system into components will help you brainstorm the hazards.
You shouldn't do a full design of the components, just get a feel for the major
ones.  For projects that involve hardware, the components will typically include
each individual piece of hardware.  If your software will have a database, or an
important library, these are also potential components.}

\section{Critical Assumptions}

\wss{These assumptions that are made about the software or system.  You should
minimize the number of assumptions that remove potential hazards.  For instance,
you could assume a part will never fail, but it is generally better to include
this potential failure mode.}

\section{Failure Mode and Effect Analysis}

\wss{The safety requirements in the table do not have to have the prefix SR.
The most important thing is to show traceability to your SRS. You might trace to
requirements you have already written, or you might need to add new
requirements.}
\wss{If no safety requirement can be devised, other mitigation strategies can be
entered in the table, including strategies involving providing additional
documentation, and/or test cases.}

\section{Safety and Security Requirements}

\noindent The requirements below are in addition to safeguards already in the SRS (see Table 4 of the SRS), including user review of extracted data (FR5), flagging of missing fields (FR4), instructor approval of date changes (FR12), login (FR15), role-based access (FR16), and restrictions on data sent to external services (AIR2). Additional safety and security requirements are listed in Table~\ref{TblSafetySecurityRequirements}. Public outline text may be sent to Copilot, but confidential data and student information shall remain internal (see DR2).

\begin{table}[H]
\caption{Safety and Security Requirements}
\label{TblSafetySecurityRequirements}
\begin{tabularx}{\textwidth}{p{0.12\textwidth}p{0.20\textwidth}X}
\toprule
\textbf{ID} & \textbf{Hazard(s)} & \textbf{Requirement}\\
\midrule
\refstepcounter{ssrnum}\ssrthessrnum \label{SSR_ProcessFailure} & Outline processing failure resulting in an invalid or incomplete schedule or lost data & The system shall display an actionable error when a course outline PDF cannot be read or the model fails to process it, shall not create assessments for that outline automatically, and shall keep all data already entered.\\
\refstepcounter{ssrnum}\ssrthessrnum \label{SSR_AcceptedFiles} & Unsupported or malicious input file & The system shall accept only PDF files smaller than 5 MB, and shall reject any other file before processing it.\\
\refstepcounter{ssrnum}\ssrthessrnum \label{SSR_RestrictedDates} & Assessment scheduled during a restricted period & The system shall use the academic calendar and shall not suggest an assessment date that falls on a holiday, recess, examination period, or test and examination ban period.\\
\refstepcounter{ssrnum}\ssrthessrnum \label{SSR_StaleSuggestion} & Out-of-date scheduling suggestion & The system shall mark a date suggestion as out of date if the assessment's date, or the other assessments in the affected weeks, change before the instructor decides on it.\\
\refstepcounter{ssrnum}\ssrthessrnum \label{SSR_CheckCoverage} & Incomplete compliance check or falsely reported absence of findings & The system shall show which compliance rules were checked for each outline, and shall state that a course was not checked when it is missing from the master course sheet, rather than reporting no findings.\\
\refstepcounter{ssrnum}\ssrthessrnum \label{SSR_AuditTrail} & Unauthorized or untraceable data modification & The system shall record the user, the time, and the previous value whenever an assessment date or extracted value is changed, and shall record each change and retain that audit record for the retention period in DR4.\\
\refstepcounter{ssrnum}\ssrthessrnum \label{SSR_MasterSheetEdit} & Unauthorized modification of course-rule data & The system shall allow only administrators to change the master course sheet.\\
\bottomrule
\end{tabularx}
\end{table}

\section{Roadmap}

\wss{Which safety requirements will be implemented as part of the capstone timeline?
Which requirements will be implemented in the future?}

\newpage{}

\section{Appendix --- Generative AI Use Disclosure}

\wss{ChatGPT (version 5) was used to brainstorm ideas for the template for this disclose section.}

\subsection{AI Tools Used}

Grok (version 4.7) was used.

\subsection{How and Where Used}

Grok was used to fix formatting issues and address pull request comments in this
Hazard Analysis document, section 6.


\subsection{Verification of AI-Generated Content}

\wss{\begin{itemize}
    \item Describe how the team checked AI-generated or AI-assisted material
    before including it in the document.
    \item In particular, verify factual claims, technical assertions,
    requirements, calculations, references, and other information for which an
    AI system may produce incorrect or fabricated results.
    \item \textbf{The team is responsible for the accuracy, completeness,
consistency, and appropriateness of the submitted document, regardless of
whether material was generated by a human or an AI system.}
\end{itemize}}

\newpage{}

\section{Appendix --- Reflection}

\wss{Not required for CAS 741}

\input{../Reflection.text}

\begin{enumerate}
    \item What went well while writing this deliverable? 
    \item What pain points did you experience during this deliverable, and how
    did you resolve them?
    \item Which of your listed risks had your team thought of before this
    deliverable, and which did you think of while doing this deliverable? For
    the latter ones (ones you thought of while doing the Hazard Analysis), how
    did they come about?
    \item Other than the risk of physical harm (some projects may not have any
    appreciable risks of this form), list at least 2 other types of risk in
    software products. Why are they important to consider?
\end{enumerate}

\end{document}